% SOURCE_STATEMENT: "Given a spatial tensor product subsystem $E^{\\odot}$ of a free product system, is $E^{\\odot}$ recovered by the intersection procedure over all decompositions of each time parameter? More generally, characterize when a free product system can be obtained in this way from a spatial tensor product subsystem."
% Preparation and final self-audit: CONTRIBUTING.md prompts C and D; source/audit artifacts retained in the local physics_statements_2026_10_11 batch.
% !TeX program = lualatex
\documentclass[11pt]{article}
\usepackage[margin=25mm]{geometry}
\usepackage{fontspec}
\setmainfont{Latin Modern Roman}
\usepackage{amsmath,amssymb,mathtools}
\usepackage{unicode-math}
\setmathfont{Latin Modern Math}
\usepackage{xurl,hyperref}
\hypersetup{colorlinks=true,urlcolor=blue}
\setcounter{secnumdepth}{0}
\setlength{\parindent}{0pt}
\setlength{\parskip}{5pt}
\setlength{\emergencystretch}{3em}
\newcommand{\sref}[2]{\hyperlink{source:#1:#2}{[S#2]}}
\newcommand{\cataloguescope}[1]{\paragraph{Recorded target.}#1}
\newcommand{\R}{\mathbb R}
\newcommand{\C}{\mathbb C}
\newcommand{\N}{\mathbb N}
\newcommand{\Z}{\mathbb Z}
\newcommand{\Q}{\mathbb Q}
\newcommand{\F}{\mathbb F}
\newcommand{\eps}{\varepsilon}
\begin{document}
% UnsolvedMath ID 30001180; source catalogue code OWR-3392-006
\setcounter{section}{30001179}
\section{Recovering Spatial Tensor Systems by Intersection}\label{30001180}
{\small UnsolvedMath ID 30001180 · Mathematical Physics}
\subsection{Definitions and mathematical statement}

\paragraph{Shared notation.}$\mathbb N=\{1,2,\ldots\}$; the standard number fields have their usual meanings. All additional conventions are stated locally.


\paragraph{Definitions and conventions.} For a unital coefficient algebra $B$, a Hilbert $B$-bimodule is a right Hilbert $B$-module with a compatible left action. $\mathcal B^a(E)$ denotes its adjointable operators. A pointed free product system $(E_t,\omega_t)$ has central reference vectors and associative identifications $E_s*E_t\simeq E_{s+t}$, using the source's reduced free product over $B$. A tensor product subsystem instead has $E_s\odot E_t\simeq E_{s+t}$ with $\odot$ the internal tensor product over $B$. An $E_0$-semigroup is a continuous semigroup of unital $*$-endomorphisms. A unit is a section whose values multiply according to the product-system structure. Starting with a spatial tensor product system $E^\odot$ and its reference unit, the source forms free products over time decompositions and takes an inductive limit to obtain a generated free product system $E^*$. Recovery asks whether intersecting all embedded tensor products at fixed total time gives exactly the original $E^\odot$, rather than a larger tensor subsystem.

\paragraph{Source scope and qualifications.} Generation by the starting system is a theorem of the construction in this report. Recovering the starting tensor system by intersection, and characterizing free systems admitting that description, are the unresolved questions. This reverses a construction, whereas ID 30001178 asks generation inside a given free system.

\paragraph{Catalogue mathematical target.}
\paragraph{Catalogue wording retained for traceability.}
Given a spatial tensor product subsystem $E^{\odot}$ of a free product system, is $E^{\odot}$ recovered by the intersection procedure over all decompositions of each time parameter? More generally, characterize when a free product system can be obtained in this way from a spatial tensor product subsystem.

\paragraph{Saved literature context.} The supplied catalogue's assessment dated 2026-08-21 records: The report gives a positive construction in a spatial setting but leaves open when the intersection procedure recovers the original tensor product system in general. The assessment is retained as dated provenance; it is not a new comprehensive literature or open-status certification.
\subsection{Short English statement}
Determine when constructing a free product system from a spatial tensor system and then intersecting tensor products recovers that original system.
\subsection{Sources}
\begin{itemize}
\item[\hypertarget{source:30001180:1}{S1}] ULAM.AI and the UnsolvedMath Contributors, UnsolvedMath, record 30001180 (OWR-3392-006). Catalogue curation under CC BY 4.0; definitions and English exposition added locally.
\url{https://huggingface.co/datasets/ulamai/UnsolvedMath}.
\item[\hypertarget{source:30001180:2}{S2}] Mini-Workshop: Product Systems and Independence in Quantum Dynamics (2009), Oberwolfach Reports, DOI 10.4171/owr/2009/09, pp. 493--548; official report source and mathematical context.
\url{https://doi.org/10.4171/owr/2009/09}.
\item[\hypertarget{source:30001180:3}{S3}] Official EMS Press PDF of the same Oberwolfach report; PDF page 38, printed page 530.
\url{https://ems.press/content/serial-article-files/46209}.
\item[\hypertarget{source:30001180:4}{S4}] free product systems; reference in the saved dated catalogue literature review.
\url{https://iris.unimol.it/handle/11695/18113}.
\end{itemize}
\label{end:30001180}
\end{document}
